\documentclass[12pt,a4paper]{report}

\usepackage[utf8]{inputenc}
\usepackage[T1]{fontenc}
\usepackage[spanish]{babel}
\usepackage{geometry}
\usepackage{setspace}
\usepackage{graphicx}
\usepackage{booktabs}
\usepackage{amsmath}
\usepackage{amssymb}
\usepackage{hyperref}
\usepackage{enumitem}

\geometry{
    left=3cm,
    right=2cm,
    top=3cm,
    bottom=2cm
}

\onehalfspacing

\begin{document}

% ============================================================
% PORTADA
% ============================================================

\begin{titlepage}
    \centering
    
    {\Large Universidad XXXXX\par}
    \vspace{0.5cm}
    
    {\Large Programa de Doctorado en XXXXX\par}
    
    \vspace{3cm}
    
    {\LARGE \textbf{Título provisional de la investigación}\par}
    
    \vspace{1cm}
    
    {\large
    Integración de información generada por Inteligencia Artificial
    en arquitecturas de control de sistemas dinámicos
    \par}
    
    \vfill
    
    {\large Nombre del doctorando\par}
    {\large Director / Orientador\par}
    
    \vspace{1cm}
    
    {\large Año\par}
    
\end{titlepage}


% ============================================================
% RESUMEN
% ============================================================

\begin{abstract}

% Escribir una síntesis muy breve de todo el proyecto.
%
% Debe contener:
%
% 1. El contexto general:
%    - uso creciente de técnicas de IA para obtener información
%      sobre el comportamiento futuro de sistemas dinámicos.
%
% 2. El problema:
%    - generar predicciones con IA no implica automáticamente
%      que estas predicciones mejoren un sistema de control.
%
% 3. La brecha de investigación:
%    - falta comprender cómo transformar e integrar la información
%      producida por IA dentro de arquitecturas de control.
%
% 4. La pregunta principal de investigación.
%
% 5. La estrategia general propuesta para responderla.
%
% 6. El tipo de sistemas o estudios de caso que serán utilizados.
%
% 7. La principal contribución científica esperada.
%
% Escribir esta sección solamente después de terminar
% el resto del documento.

\end{abstract}


\tableofcontents


% ============================================================
% CAPÍTULO 1
% ============================================================

\chapter{Introducción}

% Presentar el contexto general de la investigación.
%
% Construir el argumento desde lo general hacia lo específico:
%
% 1. Importancia de los sistemas de control en sistemas dinámicos.
%
% 2. Limitaciones de arquitecturas de control cuando:
%    - existen incertidumbres;
%    - el modelo es imperfecto;
%    - existen perturbaciones;
%    - existen dinámicas difíciles de modelar;
%    - el comportamiento futuro del sistema no es conocido.
%
% 3. Crecimiento del uso de Inteligencia Artificial y Machine Learning
%    para:
%    - predicción;
%    - estimación;
%    - identificación;
%    - detección;
%    - generación de información sobre estados futuros del sistema.
%
% 4. Explicar que gran parte de los trabajos se concentra en
%    mejorar la calidad de la predicción.
%
% 5. Introducir el problema central:
%
%    La existencia de una predicción más precisa no garantiza
%    necesariamente una mejora en el desempeño del sistema de control.
%
% 6. Plantear la necesidad de estudiar la INTERFAZ entre:
%
%       Inteligencia Artificial
%              ↓
%       información generada
%              ↓
%       arquitectura de control
%              ↓
%       decisión de control
%              ↓
%       desempeño del sistema
%
% 7. Terminar la introducción conduciendo directamente
%    hacia la pregunta de investigación.


% ------------------------------------------------------------
\section{Contextualización}

% Explicar el contexto tecnológico y científico en el que surge
% el problema.
%
% Describir:
%
% - sistemas dinámicos;
% - control automático;
% - disponibilidad creciente de datos;
% - capacidad predictiva de modelos de IA;
% - aparición de controladores que utilizan predicciones,
%   estimaciones o información aprendida.
%
% Todavía no profundizar en algoritmos específicos.


% ------------------------------------------------------------
\section{Problema de investigación}

% Definir con precisión cuál es el problema científico.
%
% Diferenciar claramente:
%
% PROBLEMA 1:
% obtener información mediante IA.
%
% PROBLEMA 2:
% utilizar correctamente esa información dentro del control.
%
% Dejar claro que la tesis se concentra principalmente
% en el segundo problema.
%
% Discutir situaciones como:
%
% - una predicción puede contener errores;
% - la incertidumbre de la predicción puede variar;
% - la información puede tener diferentes horizontes temporales;
% - una predicción útil desde el punto de vista estadístico puede
%   no ser útil desde el punto de vista del control;
% - introducir información predictiva puede modificar estabilidad,
%   robustez, seguridad o desempeño.


% ------------------------------------------------------------
\section{Brecha de investigación}

% Esta es una de las partes más importantes del examen.
%
% No escribir simplemente:
% ``pocos trabajos estudiaron este tema''.
%
% Construir la brecha a partir de la revisión de literatura.
%
% Identificar:
%
% 1. Qué sabe actualmente la literatura.
%
% 2. Qué soluciones existen para integrar IA y control.
%
% 3. Qué tipo de información de IA se utiliza actualmente.
%
% 4. Cómo se introduce esa información en el controlador.
%
% 5. Qué criterios se utilizan para decidir si la información
%    realmente mejora el control.
%
% 6. Qué aspectos siguen sin estar suficientemente explicados.
%
% La brecha debería conducir al problema:
%
% Falta un entendimiento/metodología/arquitectura sistemática
% para determinar cómo la información generada por modelos de IA
% debe ser incorporada al sistema de control para producir
% mejoras verificables de desempeño.


% ------------------------------------------------------------
\section{Pregunta de investigación}

% Escribir aquí la pregunta principal.
%
% Versión preliminar:
%
% ¿Cómo utilizar la información generada por modelos de
% Inteligencia Artificial para mejorar el desempeño
% de sistemas de control?
%
% Posteriormente, refinar la formulación.
%
% Una versión más específica podría incluir:
%
% - integración;
% - arquitectura de control;
% - sistemas dinámicos;
% - desempeño;
% - incertidumbre de las predicciones.


% ------------------------------------------------------------
\section{Preguntas específicas de investigación}

% Descomponer la pregunta principal.
%
% No es necesario que estas sean exactamente las preguntas finales.
%
% Investigar preguntas como:
%
% P1:
% ¿Qué tipos de información generada por IA pueden ser relevantes
% para un sistema de control?
%
% P2:
% ¿En qué parte de una arquitectura de control puede integrarse
% esa información?
%
% P3:
% ¿Cómo afectan la precisión, incertidumbre y horizonte de predicción
% al desempeño del controlador?
%
% P4:
% ¿Bajo qué condiciones la información generada por IA produce
% una mejora real del sistema de control?
%
% P5:
% ¿Cómo comparar una arquitectura convencional con una arquitectura
% asistida por información generada mediante IA?


% ------------------------------------------------------------
\section{Hipótesis o proposición de investigación}

% Definir la hipótesis central que será evaluada.
%
% Estructura sugerida:
%
% Si información relevante sobre el comportamiento futuro
% o estados no observables del sistema es obtenida mediante IA
% y se integra adecuadamente en la arquitectura de control,
% entonces es posible mejorar determinados indicadores de desempeño
% respecto a una arquitectura que no utiliza dicha información.
%
% Después especificar:
%
% - qué significa ``información relevante'';
% - qué significa ``integrada adecuadamente'';
% - qué significa ``mejorar'';
% - cuáles indicadores medirán esa mejora.


% ============================================================
% CAPÍTULO 2
% ============================================================

\chapter{Objetivos}


\section{Objetivo general}

% Escribir un único objetivo general.
%
% El objetivo debe responder directamente a la pregunta
% de investigación.
%
% Debe contener tres elementos:
%
% 1. Acción científica:
%    desarrollar / investigar / proponer / evaluar.
%
% 2. Objeto:
%    integración de información generada mediante IA.
%
% 3. Finalidad:
%    mejorar el desempeño de sistemas de control.


\section{Objetivos específicos}

% Crear objetivos específicos que representen etapas necesarias
% para alcanzar el objetivo general.
%
% Una posible estructura conceptual:
%
% OE1:
% caracterizar los tipos de información generada por IA
% utilizados en sistemas de control.
%
% OE2:
% identificar y clasificar las formas de integración entre
% modelos de IA y arquitecturas de control.
%
% OE3:
% definir una arquitectura o metodología para incorporar
% información predictiva al sistema de control.
%
% OE4:
% analizar cómo propiedades de la información generada por IA
% afectan el desempeño del control.
%
% OE5:
% implementar la propuesta en sistemas dinámicos seleccionados.
%
% OE6:
% comparar experimentalmente/simuladamente el desempeño con
% arquitecturas de referencia.


% ============================================================
% CAPÍTULO 3
% ============================================================

\chapter{Fundamentación teórica}

% Este capítulo debe definir los conceptos necesarios
% para entender el problema.
%
% No debe convertirse todavía en una revisión exhaustiva
% de todos los artículos.
%
% Cada sección debe terminar mostrando cómo ese concepto
% se relaciona con la pregunta de investigación.


% ------------------------------------------------------------
\section{Sistemas dinámicos}

% Definir qué se entiende por sistema dinámico en esta tesis.
%
% Presentar una representación matemática general.
%
% Explicar:
%
% - estados;
% - entradas;
% - salidas;
% - perturbaciones;
% - incertidumbres.
%
% Definir posteriormente qué clase de sistemas será considerada.


% ------------------------------------------------------------
\section{Sistemas de control}

% Presentar los conceptos fundamentales necesarios.
%
% Incluir:
%
% - lazo abierto y cerrado;
% - realimentación;
% - referencia;
% - controlador;
% - planta;
% - perturbaciones;
% - estabilidad;
% - seguimiento;
% - robustez;
% - desempeño.
%
% No realizar una revisión enciclopédica.
% Presentar solamente los conceptos necesarios para la tesis.


% ------------------------------------------------------------
\section{Arquitecturas de control}

% Definir qué significa ``arquitectura de control''
% dentro del trabajo.
%
% Explicar dónde pueden aparecer diferentes módulos:
%
% sensores
%      ↓
% estimación
%      ↓
% predicción
%      ↓
% controlador
%      ↓
% actuadores
%      ↓
% planta
%
% Discutir diferentes puntos posibles de integración de IA.


% ------------------------------------------------------------
\section{Inteligencia Artificial aplicada a sistemas dinámicos}

% Delimitar qué se entiende por IA en la tesis.
%
% Evitar una historia general de la Inteligencia Artificial.
%
% Concentrarse en modelos capaces de generar información útil
% para el control, por ejemplo:
%
% - Machine Learning;
% - Deep Learning;
% - redes neuronales;
% - modelos data-driven;
% - modelos híbridos;
% - modelos predictivos.
%
% Explicar qué función cumplen dentro del problema estudiado.


% ------------------------------------------------------------
\section{Predicción mediante Inteligencia Artificial}

% Definir formalmente qué se entiende por predicción.
%
% Especificar posibles:
%
% - variables predichas;
% - horizonte de predicción;
% - predicción de estados;
% - predicción de perturbaciones;
% - predicción de referencias;
% - predicción de comportamiento futuro.
%
% Diferenciar:
%
% calidad de predicción
%
% de
%
% utilidad de la predicción para el control.


% ------------------------------------------------------------
\section{Información generada por Inteligencia Artificial}

% Esta sección es especialmente importante para la tesis.
%
% Definir qué significa exactamente ``información generada
% por la IA''.
%
% Construir una taxonomía preliminar.
%
% Por ejemplo:
%
% - estados estimados;
% - estados futuros;
% - perturbaciones futuras;
% - parámetros estimados;
% - clasificación de condiciones de operación;
% - detección de anomalías;
% - probabilidades;
% - incertidumbre;
% - trayectorias futuras;
% - restricciones estimadas.
%
% Esta clasificación deberá estar fundamentada posteriormente
% mediante la revisión de literatura.


% ============================================================
% CAPÍTULO 4
% ============================================================

\chapter{Estado del arte}

% Aquí realizar la revisión crítica de la literatura.
%
% No organizar únicamente por autor o cronológicamente.
%
% Organizar los trabajos según los conceptos relevantes
% para responder la pregunta de investigación.


% ------------------------------------------------------------
\section{Estrategia de revisión de literatura}

% Explicar cómo se realizó la búsqueda bibliográfica.
%
% Informar:
%
% - bases de datos;
% - período considerado;
% - palabras clave;
% - strings de búsqueda;
% - criterios de inclusión;
% - criterios de exclusión;
% - proceso de selección.
%
% Mostrar posteriormente un diagrama del proceso si corresponde.


% ------------------------------------------------------------
\section{IA para predicción en sistemas dinámicos}

% Analizar trabajos que utilizan IA para predecir variables
% relacionadas con sistemas dinámicos.
%
% Para cada grupo de trabajos identificar:
%
% - qué se predice;
% - modelo de IA;
% - horizonte;
% - precisión;
% - aplicación;
% - si la predicción entra o no al sistema de control.


% ------------------------------------------------------------
\section{Integración entre IA y sistemas de control}

% Analizar específicamente cómo la salida del modelo de IA
% se conecta con el sistema de control.
%
% Clasificar las diferentes arquitecturas encontradas.
%
% Por ejemplo:
%
% IA como estimador;
% IA como predictor;
% IA como modelo de planta;
% IA para anticipar perturbaciones;
% IA para ajustar parámetros del controlador;
% IA integrada directamente en la política de control.


% ------------------------------------------------------------
\section{Uso de información predictiva en control}

% Analizar trabajos donde el controlador utiliza conocimiento
% sobre eventos o estados futuros.
%
% Comparar especialmente con:
%
% - control predictivo;
% - feedforward;
% - disturbance preview;
% - preview control;
% - anticipatory control;
% - predictive control basado en modelos aprendidos.
%
% Identificar similitudes y diferencias con la propuesta doctoral.


% ------------------------------------------------------------
\section{Evaluación del impacto de la IA sobre el control}

% Investigar cómo los trabajos existentes demuestran
% que utilizar IA mejora el sistema.
%
% Registrar indicadores utilizados:
%
% - error de seguimiento;
% - IAE;
% - ISE;
% - RMSE;
% - tiempo de establecimiento;
% - overshoot;
% - consumo de energía;
% - esfuerzo de control;
% - estabilidad;
% - robustez;
% - seguridad;
% - cumplimiento de restricciones.
%
% Analizar si los artículos separan:
%
% desempeño predictivo
%
% de
%
% desempeño de control.


% ------------------------------------------------------------
\section{Síntesis crítica del estado del arte}

% Construir una tabla comparativa de los principales trabajos.
%
% Columnas sugeridas:
%
% Referencia
% Sistema
% Técnica de IA
% Información generada
% Horizonte
% Arquitectura de control
% Forma de integración
% Métrica de predicción
% Métrica de control
% Comparación con baseline
% Limitaciones
%
% A partir de esta tabla construir la brecha de investigación.


% ============================================================
% CAPÍTULO 5
% ============================================================

\chapter{Marco conceptual propuesto}

% Presentar aquí la forma en que la tesis organiza conceptualmente
% el problema.
%
% Este capítulo debe conectar todos los conceptos anteriores.


\section{Cadena información--decisión--control}

% Proponer un esquema conceptual como:
%
% Sistema dinámico
%       ↓
% Datos
%       ↓
% Modelo de IA
%       ↓
% Información generada
%       ↓
% Mecanismo de integración
%       ↓
% Controlador
%       ↓
% Acción de control
%       ↓
% Desempeño del sistema
%
% Explicar cada bloque.


\section{Caracterización de la información generada por IA}

% Proponer dimensiones para caracterizar la información.
%
% Por ejemplo:
%
% - tipo de variable;
% - horizonte temporal;
% - frecuencia de actualización;
% - error;
% - incertidumbre;
% - confiabilidad;
% - relevancia para el controlador.


\section{Mecanismos de integración}

% Definir las formas posibles en las que la información
% puede modificar la decisión de control.
%
% Ejemplos conceptuales:
%
% - modificar una referencia;
% - modificar una restricción;
% - modificar un parámetro del controlador;
% - proporcionar una perturbación futura;
% - proporcionar estados futuros;
% - modificar una función de costo;
% - activar diferentes estrategias de control.


% ============================================================
% CAPÍTULO 6
% ============================================================

\chapter{Metodología de investigación}

% Explicar cómo se responderá científicamente
% la pregunta de investigación.
%
% Esta sección debe permitir que la banca evalúe
% si realmente es posible verificar la hipótesis.


\section{Estrategia general}

% Presentar las etapas principales de la investigación.
%
% Ejemplo:
%
% Etapa 1:
% revisión sistemática y construcción de taxonomía.
%
% Etapa 2:
% definición de la arquitectura conceptual.
%
% Etapa 3:
% selección/modelado de sistemas dinámicos.
%
% Etapa 4:
% desarrollo de modelos de IA.
%
% Etapa 5:
% integración con el sistema de control.
%
% Etapa 6:
% diseño de experimentos.
%
% Etapa 7:
% evaluación y comparación.


\section{Sistemas dinámicos de estudio}

% Especificar qué sistemas serán utilizados para validar
% la propuesta.
%
% Para cada sistema:
%
% - descripción;
% - estados;
% - entradas;
% - salidas;
% - perturbaciones;
% - modelo;
% - restricciones;
% - por qué es relevante para la tesis.
%
% Justificar la elección de los casos de estudio.


\section{Arquitectura de control de referencia}

% Definir el baseline.
%
% Explicar cómo funcionará el sistema SIN utilizar
% información generada mediante IA.
%
% Esta arquitectura será fundamental para demostrar
% posteriormente si existe mejora.


\section{Modelo de Inteligencia Artificial}

% Describir qué información será generada.
%
% No concentrarse solamente en qué algoritmo tiene
% la menor RMSE.
%
% Informar:
%
% - entradas del modelo;
% - salida;
% - datos de entrenamiento;
% - validación;
% - horizonte;
% - incertidumbre;
% - métricas de desempeño predictivo.


\section{Arquitectura de control asistida por IA}

% Describir la arquitectura propuesta.
%
% Indicar exactamente:
%
% 1. qué información produce la IA;
%
% 2. dónde entra esa información;
%
% 3. cómo modifica la decisión de control;
%
% 4. qué ocurre cuando la predicción es incorrecta;
%
% 5. qué ocurre si no existe información disponible;
%
% 6. qué mecanismos protegen estabilidad/robustez/seguridad.


\section{Diseño experimental}

% Definir los escenarios que serán comparados.
%
% Como mínimo considerar:
%
% C0:
% controlador convencional.
%
% C1:
% controlador + información ideal/perfecta.
%
% C2:
% controlador + predicción de IA.
%
% Esto permite separar:
%
% potencial máximo de la información
%
% de
%
% desempeño real del modelo de IA.
%
% También considerar diferentes niveles de:
%
% - error de predicción;
% - horizonte;
% - ruido;
% - incertidumbre;
% - perturbaciones.


\section{Métricas de evaluación}

% Definir previamente qué significa ``mejorar el control''.
%
% Seleccionar métricas apropiadas para los sistemas estudiados.
%
% Separar explícitamente:
%
% MÉTRICAS DE IA
%
% - RMSE;
% - MAE;
% - R²;
% - accuracy;
% - otras.
%
% de
%
% MÉTRICAS DE CONTROL
%
% - error de seguimiento;
% - IAE;
% - ISE;
% - overshoot;
% - tiempo de establecimiento;
% - consumo energético;
% - esfuerzo de control;
% - robustez;
% - violación de restricciones.


\section{Análisis de sensibilidad}

% Investigar cómo la calidad de la información modifica
% el resultado del sistema de control.
%
% Variar sistemáticamente:
%
% - error;
% - horizonte de predicción;
% - frecuencia de actualización;
% - incertidumbre;
% - retraso;
% - ruido.
%
% Buscar identificar regiones donde:
%
% IA mejora el control;
%
% IA no produce diferencia;
%
% IA perjudica el control.


% ============================================================
% CAPÍTULO 7
% ============================================================

\chapter{Resultados preliminares}

% Incluir solamente resultados ya obtenidos hasta el examen.
%
% Mostrar resultados que ayuden a demostrar la viabilidad
% de la investigación.
%
% Pueden incluir:
%
% - implementación inicial del sistema;
% - modelo dinámico;
% - controlador baseline;
% - primeros modelos de IA;
% - desempeño predictivo;
% - integración preliminar;
% - comparación inicial.
%
% Evitar presentar resultados aislados sin relacionarlos
% con la pregunta de investigación.


\section{Discusión de los resultados preliminares}

% Explicar qué demuestran realmente los resultados obtenidos.
%
% Separar:
%
% lo que ya puede afirmarse;
%
% lo que todavía es una hipótesis;
%
% lo que todavía debe ser investigado.


% ============================================================
% CAPÍTULO 8
% ============================================================

\chapter{Contribuciones esperadas}

% Explicar cuál será el conocimiento nuevo generado por la tesis.
%
% Evitar presentar como contribución solamente:
%
% ``desarrollar un modelo de IA''.
%
% Las contribuciones deberían estar relacionadas con
% la pregunta de investigación.
%
% Posibles tipos de contribución:
%
% 1. Taxonomía de información generada por IA relevante para control.
%
% 2. Marco conceptual para integrar esa información
%    en arquitecturas de control.
%
% 3. Metodología para evaluar cuándo una predicción
%    es realmente útil para control.
%
% 4. Arquitectura de control que utilice información predictiva.
%
% 5. Relación entre calidad de predicción y desempeño de control.
%
% 6. Criterios para determinar cuándo la utilización de IA
%    mejora, no altera o perjudica un sistema de control.


% ============================================================
% CAPÍTULO 9
% ============================================================

\chapter{Plan de trabajo}

% Presentar las actividades que todavía deben ser realizadas
% después del examen de calificación.
%
% Relacionar cada actividad con un objetivo específico.


\section{Actividades}

% Crear una lista/tablas con:
%
% - actividad;
% - objetivo asociado;
% - resultado esperado;
% - período.


\section{Cronograma}

% Insertar cronograma del tiempo restante del doctorado.
%
% Incluir:
%
% - revisión bibliográfica;
% - desarrollo teórico;
% - implementación;
% - experimentos;
% - análisis;
% - publicaciones;
% - escritura de tesis;
% - defensa.


% ============================================================
% CAPÍTULO 10
% ============================================================

\chapter{Riesgos y limitaciones}

% Identificar riesgos científicos y metodológicos.
%
% Por ejemplo:
%
% - insuficiencia de datos;
% - generalización limitada del modelo de IA;
% - errores de predicción;
% - dependencia del sistema estudiado;
% - dificultad de demostrar estabilidad;
% - costo computacional;
% - problemas de tiempo real.
%
% Para cada riesgo explicar brevemente
% cómo será mitigado.


% ============================================================
% CAPÍTULO 11
% ============================================================

\chapter{Conclusiones del examen de calificación}

% No escribir las conclusiones finales de la tesis.
%
% Resumir:
%
% 1. cuál es el problema científico;
%
% 2. cuál es la brecha identificada;
%
% 3. cuál es la pregunta de investigación;
%
% 4. cuál es la hipótesis;
%
% 5. cómo se pretende responderla;
%
% 6. qué evidencia preliminar existe;
%
% 7. qué falta realizar;
%
% 8. cuál será la contribución esperada.


% ============================================================
% REFERENCIAS
% ============================================================

\bibliographystyle{plain}
\bibliography{referencias}

\end{document}
